\begingroup
\setlength{\LTcapwidth}{\textwidth}
\begin{longtable}{p{3.2cm}p{4.2cm}p{5.0cm}}
\caption{Kerangka evaluasi komponen prediksi dan RAG dalam sistem pendukung keputusan.}
\label{tab:bab2.4-kerangka-evaluasi} \\
\toprule
\textbf{Komponen}
        & \textbf{Metrik/Instrumen}
        & \textbf{Aspek yang Diukur} \\
\midrule
\endfirsthead

\multicolumn{3}{@{}l}{\small\emph{Tabel \thetable{} (lanjutan)}} \\
\toprule
\textbf{Komponen}
        & \textbf{Metrik/Instrumen}
        & \textbf{Aspek yang Diukur} \\
\midrule
\endhead

\bottomrule
\endlastfoot

Prediksi glukosa
&
RMSE, MAE, MAPE
&
Besarnya kesalahan numerik antara nilai prediksi dan nilai aktual.
\\

Prediksi klinis
&
Clarke Error Grid
&
Konsekuensi klinis dari perbedaan antara nilai prediksi dan nilai aktual.
\\

Retrieval
&
Hit@k, MRR, nDCG
&
Keberadaan, posisi, dan kualitas pemeringkatan dokumen relevan.
\\

Generation
&
Faithfulness, context precision, context recall
&
Keterikatan keluaran terhadap konteks serta kualitas konteks yang digunakan.
\\

Stability
&
Repeated-run stability
&
Konsistensi keluaran sistem pada eksekusi berulang.
\\

Safety
&
Generation safety
&
Kemunculan keluaran yang tidak didukung, melampaui konteks, atau tidak menyatakan keterbatasan.
\\

Usability
&
System Usability Scale (SUS)
&
Persepsi kemudahan penggunaan sistem oleh pengguna sasaran.
\\
\end{longtable}
\endgroup
